\scanpage{34}\mspage{33}
% L1–3
\deleted{\unsure{7|8}}8) $\forall\,X\in\mathcal M$, soit $X'$ son ``complémentaire'' \deleted{\unsure{8|6}} Comme $\{X,X'\}$ modéré, alors $\{\overline X,\overline{X'}\}$ \deleted{$\mathrm{Fr}\,X$ est \unsure{fini} $\unsure{\{\overline X,\overline{X'}\}|\mathrm{Fr}\,X}\in\mathcal M$} \inserted{aussi, on peut} former
\[
\left\{\begin{aligned}
\mathrm{Fr}\,X &\overset{\text{def}}{=} \overline X\wedge\overline{X'}\\
(\mathrm{Fr}\,X)' &= X^{\circ}\vee X^{\prime\circ}
\end{aligned}\right.
\]
% slanted
\margin{propriété spéciale au cas envisagé du modèle de \unsure{formes|fermés} \uline{de} \uline{dim $1$}}
% L4–5
Ceci dit $\{X,\deleted{\unsure{X'}}\unsure{Y|Y'}\}$ modéré \deleted{\unsure{ssi}} $\{\mathrm{Fr}\,X,\mathrm{Fr}\,Y\}$ est modéré $\Longrightarrow$ $\mathrm{Fr}\,X\cup\mathrm{Fr}\,X$ est modéré.

% L6–10
\deleted{\uline{NB}} \deleted{Soit $X$ \unsure{fermé|modéré}, alors} \deleted{\inserted{On dit que $X$ \unsure{($\in\mathcal M$)|(de $\mathcal M$)} est \uline{rare} si}} \deleted{si $X^\circ=\varnothing$ i.e. $\overline{X'}=1$ ($X'$ \uline{dense}). Dans le cas d'espèce, $X$ est \unsure{rare|un} sss $X$ est fini (en tant que partie de $L$), ou encore sss $\mathrm{Fr}\,X$}

% L11–12
Dans le cas d'espèce, conditions équivalentes sur $X\in\mathcal M$
% L13 (a)
\[
\begin{aligned}
\text{a)}\quad & X\subset\deleted{\unsure{\mathrm{Fr}(X)}}X^{\bullet}
&&\left(\begin{aligned}&\text{i.e.}\quad X^{\circ}=\varnothing\quad\&\ \text{i.e.}\quad\overline{X'}=1\\[-2pt]&\hspace{8em}\text{i.e. }X'\text{ ``dense''}\end{aligned}\right)\\
% L14 (b)
\text{b)}\quad & X=X^{\bullet}
&&\left(\begin{aligned}&\text{i.e. }X\text{ fermé rare \quad i.e. }X'\text{ ouvert}\\[-2pt]&\hspace{12em}\text{dense}\end{aligned}\right)
% L15 (c)
\end{aligned}
\]
c) $X$ est fini en tant que partie de $L$

\begin{center}\rule{1.6cm}{0.45pt}\end{center}

% L16–20
Considérons un ens ordonné $\mathcal M$ satisfaisant les conditions 1) à 7) (avec 4) renforcé marginalement comme il a été dit~; dans la notation algébrique, on a
% L21–24
\diagram{\begin{tikzcd}[ampersand replacement=\&,row sep=0.1ex,column sep=1.8em,
  cells={nodes={anchor=base west}},execute at end picture={
  \draw ([xshift=-0.9em]alJoin.north west) -- ([xshift=-0.9em]alOrder.south west);
  \draw[->] ([xshift=0.1em]alOrder.north west)
      .. controls ([xshift=-0.45em]alMeet.west) and ([xshift=-0.7em]alJoin.west)
      .. ([xshift=-0.45em]alJoin.north west)
      .. controls (alJoin.north west) and ([xshift=0.5em]alJoin.north west)
      .. ([xshift=1.1em,yshift=0.25em]alJoin.north west);
  \draw ([xshift=2.6em]alOrder.north west) to[bend right=8] (alOrder.north west);
  \draw (alOrder.south west) to[bend right=8] ([xshift=2.6em]alOrder.south west);
}]
X+Y=(X\vee Y)\wedge(X\wedge Y)' \& |[alias=alJoin]| X\vee Y=X+Y-XY\deleted{\unsure{Y|X}}\\
X\cdot Y=X\wedge Y \& |[alias=alMeet]| X\wedge Y=X\cdot Y\\
0=\text{plus petit él.} \& |[alias=alOrder]| X\leq Y\Longleftrightarrow X\cdot Y=X\\
1=\text{plus gr\textsuperscript{d} él} \&
\end{tikzcd}}
